\section{Results}
\label{sec:results}
% About 3 pages. Floats hold generated bodies only; captions are prose. Every number is a generated macro.

\subsection{RQ1: Per-Document Choice versus One Fixed Setup}
\label{sec:rq1}

\begin{table*}[t]
\centering
\generated{tables/rq1.tex}
\caption{RQ1 on the held-out documents (\RQoneDocuments{} documents and \RQoneQuestions{} questions in every choice set; hybrid retrieval, Recall@1,250). Short names: Appendix~\ref{app:key}.}
\label{tab:rq1}
\end{table*}

Choosing a setup per document does not beat one fixed setup on the
held-out documents (Table~\ref{tab:rq1}). In the primary choice set of
\RQoneSetups{} setups, the best single setup is \RQoneBest{}, with
Recall@1,250 of \RQoneBestRecall{}; \RQoneTied{} of the \RQoneSetups{}
setups are tied with it. The split-half oracle reaches \RQoneOracle{},
so the gap is \RQonePrimaryGap{} \RQonePrimaryGapCI{}. The interval
contains zero, and the rule frozen before the held-out set was scored
gives \RQoneDecision{}: the recommendation is one fixed setup. Choosing
and scoring on the same questions gives a gap of \RQoneRawGap{}; this
number is inflated by selection and is shown for reference only.

On the three grids we evaluated, the gap was more negative for the
larger choice sets: \RQoneOriginalGap{} \RQoneOriginalGapCI{} on the
original grid and \RQoneFullGap{} \RQoneFullGapCI{} on the full grid of
\HeldOutSetups{} setups. A per-document choice rests on a few questions
per document, and with more setups to choose from, the choice on one
half of a document's questions holds less often on the other half.

The setup chosen on the tuning documents, \RQoneTuningWinner{}, applied
unchanged, scores \RQoneTuningWinnerRecall{} on the held-out documents:
rank \RQoneTuningWinnerRank{} of \RQoneSetups{}, \RQoneTuningWinnerTied{}
with the best. The decision is the same with BM25-only and dense-only
retrieval (Appendix~\ref{app:retrievers}) and without the Indian
documents (Appendix~\ref{app:no-indian}). The latter matters because a
check of the Indian questions found \GoldenInvalid{} of
\GoldenSampled{} invalid (Appendix~\ref{app:golden}).

\subsection{RQ2: Chunk Size and Chunker at a Fixed Budget}
\label{sec:rq2}

\begin{figure}[t]
\centering
\generated{figures/rq2.tex}
\caption{RQ2: Recall@1,250 against measured mean chunk size, one panel per parser, one mark shape per source (tuning documents).}
\label{fig:rq2}
\end{figure}

At a fixed token budget, smaller chunks retrieve more of the gold
(Figure~\ref{fig:rq2}). On the tuning documents, every source and parser
has a negative Spearman correlation between measured chunk size and
recall at every budget (\RQtwoRhoRange{}), and the best settings measure
\RQtwoBestSizes{} tokens. Held-out OHR-Bench shows the same pattern
(\RQtwoHeldOutOHRRhoRange{}; best sizes \RQtwoHeldOutOHRBestSizes{}). On
the held-out CUAD documents the correlation weakens as the budget grows;
the other held-out sources have too few documents to show a trend.

The chunker matters less than its size: in the primary choice set,
setups of different chunkers are tied with the best
(Table~\ref{tab:rq1}). At equal chunk size, Method~1 makes more chunks
than recursive splitting from the same number of index tokens, so its
gain is recall, not storage (Appendix~\ref{app:index-cost}).

\subsection{RQ3: The Measure Decides the Best Size}
\label{sec:rq3}

\begin{figure}[t]
\centering
\generated{figures/rq3.tex}
\caption{RQ3: mean chunk size of the best setting under each measure (tuning documents, settings in scope). Left of the dashed line: measures at a fixed number of chunks; right: recall at an equal budget of \FixedKBudget{} tokens. Circles: Docling; squares: pdfplumber.}
\label{fig:rq3}
\end{figure}

Which size looks best depends on the measure (Figure~\ref{fig:rq3}).
With a fixed number of chunks, larger chunks read more text, so the
fixed-$k$ measures favour them. For Recall@\TopKSmall{} the best setting
is \RQthreeDoclingRecallFiveBest{} on Docling
(\RQthreeDoclingRecallFiveScore{}) and \RQthreePdfplumberRecallFiveBest{}
on pdfplumber (\RQthreePdfplumberRecallFiveScore{}); for
Recall@\TopKLarge{} it is \RQthreeDoclingRecallOneZeroBest{} and
\RQthreePdfplumberRecallOneZeroBest{}. Hit@1 and MRR barely separate the
sizes. At an equal budget of \FixedKBudget{} tokens the ranking reverses:
the best settings are \RQthreeDoclingRecallFourZeroZeroZeroBest{} on
Docling and \RQthreePdfplumberRecallFourZeroZeroZeroBest{} on
pdfplumber, whose chunks measure \RQthreeDoclingRecallFourZeroZeroZeroSize{}
and \RQthreePdfplumberRecallFourZeroZeroZeroSize{} tokens on average.
This analysis is descriptive, on the tuning documents, and has no
intervals. It covers \RQthreeSettingsKept{} settings; the
\RQthreeSettingsLeftOut{} superseded and sanity settings are left out,
as in the main analysis (Appendix~\ref{app:erratum}).

\subsection{RQ4: Answer Quality}
\label{sec:rq4}

\begin{table*}[t]
\centering
\generated{tables/rq4.tex}
\caption{RQ4: answer token F1 on the held-out documents, each setup against the best setup, \RQfourBest{} (F1 \RQfourBestFOne{}); intervals and Holm $p$ over \RQfourComparisons{} comparisons.}
\label{tab:rq4}
\end{table*}

Answers confirm the penalty for large chunks but do not separate the
small-chunk setups (Table~\ref{tab:rq4}). On the held-out documents the
best setup by token F1 is \RQfourBest{} (\RQfourBestFOne{}). Against it,
\RQfourDoclingFixedTokenSizeOneZeroTwoFourSetup{} scores
\RQfourDoclingFixedTokenSizeOneZeroTwoFourDiff{} lower
\RQfourDoclingFixedTokenSizeOneZeroTwoFourCI{} (Holm $p$
\RQfourDoclingFixedTokenSizeOneZeroTwoFourHolm{} over
\RQfourComparisons{} comparisons).
\RQfourDoclingMethodOneTreeTThreeEightFourSetup{} is also lower, by
\RQfourDoclingMethodOneTreeTThreeEightFourDiff{}
\RQfourDoclingMethodOneTreeTThreeEightFourCI{} (Holm $p$
\RQfourDoclingMethodOneTreeTThreeEightFourHolm{}). This runs against its
recall, where it is tied with the best, and we report it without
claiming it. LangChain's library default is lower by
\RQfourDoclingLangchainDefaultDiff{} \RQfourDoclingLangchainDefaultCI{},
which is not significant after Holm ($p$
\RQfourDoclingLangchainDefaultHolm{}), although its recall loss is
(\S\ref{sec:rq5}). On the tuning documents no difference is significant
(smallest Holm $p$ \RQfourTuningMinHolm{}).

Across setups, token F1 and Recall@1,250 have a Spearman correlation of
\RQfourSpearmanTuningFOne{} on the tuning and
\RQfourSpearmanHeldOutFOne{} on the held-out documents; with the judge
as the answer measure (secondary) the values are
\RQfourSpearmanTuningJudge{} and \RQfourSpearmanHeldOutJudge{}. On the
\AQScannedQuestions{} scanned questions, the best setup,
\RQfourScannedBest{} (F1 \RQfourScannedBestFOne{}), scores above
\RQfourScannedAboveDoclingOcrTesseract{} of the
\RQfourScannedCountDoclingOcrTesseract{} Tesseract setups, but only
\RQfourScannedSignificant{} of the \RQfourScannedComparisons{}
differences is significant after Holm (\RQfourScannedTopSetup{},
\RQfourScannedTopDiff{} \RQfourScannedTopCI{}, Holm $p$
\RQfourScannedTopHolm{}). With this few questions the evidence is weak.

\subsection{RQ5: Tuned Settings versus Defaults}
\label{sec:rq5}

\begin{figure}[t]
\centering
\generated{figures/rq5.tex}
\caption{RQ5 on the held-out documents: recall difference of the tuned setting against the default with \CILevel{}\% intervals; (a) parser settings on the digital sample, (b) tuned chunk size. Filled circles: significant after Holm correction; open circles: not significant. Short names: Appendix~\ref{app:key}.}
\label{fig:rq5}
\end{figure}

\paragraph{(b) Chunk size.} LangChain's library default loses clearly
against the size tuned on the tuning documents (Figure~\ref{fig:rq5}).
On Docling, held-out recall rises from
\RQfiveBLangchainRecursiveDoclingDefault{} to
\RQfiveBLangchainRecursiveDoclingTuned{}
(\RQfiveBLangchainRecursiveDoclingDiff{}
\RQfiveBLangchainRecursiveDoclingCI{}); on pdfplumber, from
\RQfiveBLangchainRecursivePdfplumberDefault{} to
\RQfiveBLangchainRecursivePdfplumberTuned{}
(\RQfiveBLangchainRecursivePdfplumberDiff{}
\RQfiveBLangchainRecursivePdfplumberCI{}). Both are significant after
Holm over \RQfiveBComparisons{} comparisons (Holm $p$
\RQfiveBLangchainRecursiveDoclingHolm{} and
\RQfiveBLangchainRecursivePdfplumberHolm{}). For our own chunkers, the
tuned size against the starting size makes no significant difference
after Holm; the closest are Method~1 on pdfplumber
(\RQfiveBMethodOneTreePdfplumberDiff{}
\RQfiveBMethodOneTreePdfplumberCI{}) and fixed-size windows on
pdfplumber (\RQfiveBFixedTokenPdfplumberDiff{}
\RQfiveBFixedTokenPdfplumberCI{}).

\paragraph{(a) Parser settings.} No parser setting beats its default on
the digital sample: every difference is at most
\RQfiveADigitalMaxAbsDiff{} in size, and every Holm $p$ is at least
\RQfiveADigitalMinHolm{} (\RQfiveADigitalQuestions{} shared questions per
comparison). The parser itself matters more: the Docling default reaches
\RQfiveADigitalDoclingRecall{} against \RQfiveADigitalPdfplumberRecall{}
for pdfplumber on this sample. On scanned PDFs, Tesseract OCR and
Docling's default OCR do not differ (\RQfiveAScannedDiff{}
\RQfiveAScannedCI{}, Holm $p$ \RQfiveAScannedHolm{};
\RQfiveAScannedDocuments{} PDFs, \RQfiveAScannedQuestions{} shared
questions). pdfplumber reads no text from scans, so the decisive choice
is detecting a scan, not tuning the OCR; gold placement on the scans is
in Appendix~\ref{app:scanned}.

\subsection{Gold Coverage}
\label{sec:coverage}

All recall numbers are computed on questions whose gold span is located
in a parser's output. On the tuning documents Docling loses
\CoverageTuningDoclingLost{} of \CoverageTuningDoclingTotal{} main
questions (\CoverageTuningDoclingShare{}) and pdfplumber
\CoverageTuningPdfplumberLost{} (\CoverageTuningPdfplumberShare{}); on
the held-out documents Docling loses \CoverageHeldOutDoclingLost{} of
\CoverageHeldOutDoclingTotal{} (\CoverageHeldOutDoclingShare{}) and
pdfplumber \CoverageHeldOutPdfplumberLost{}
(\CoverageHeldOutPdfplumberShare{}). The held-out gap comes mostly from
the Indian documents (Docling \CoverageBySourceHeldOutIndianDocling{}
against pdfplumber \CoverageBySourceHeldOutIndianPdfplumber{}) and from
the OHR-Bench documents of the second held-out group
(\CoverageBySourceHeldOutCOHRBenchDocling{} against
\CoverageBySourceHeldOutCOHRBenchPdfplumber{}). Comparisons on shared
questions therefore hide Docling's coverage advantage on the held-out
documents.
